\chapter{PPO Formulation and Version Evolution}

Nebula's PPO subsystem evolved through three versioned formulations with
explicit compatibility checks, so changes to the state, reward, and design space
could be evaluated without silently changing historical checkpoints. All versions use separate two-layer, 64-unit tanh policy
and value networks, generalized advantage estimation with $\gamma=0.99$ and
$\lambda=0.95$, PPO clipping $\epsilon=0.2$, and entropy coefficient 0.01
\citep{schulman2017ppo,schulman2016gae}. Each parameter has an independent
categorical action head whose default choices map to grid-index deltas
$\{-1,0,+2\}$, clipped at the grid boundaries.

\section{PPO v1: AutoCkt-style baseline}

PPO v1 established the first end-to-end RL-to-SPICE path. Its five variables are
$R_{\mathrm{LOAD}}$, $R_{\mathrm{DEG}}$, $C_{\mathrm{DEG}}$,
$I_{\mathrm{TAIL}}$, and the behavioral DFE tap. Each uses a 21-point grid. An
action updates parameter index $i$ according to
\begin{equation}
  k_{i,t+1}=\operatorname{clip}\!\left(k_{i,t}+\Delta_{i,t},0,N_i-1\right),
  \qquad \Delta_{i,t}\in\{-1,0,+2\}.
\end{equation}

The 13-dimensional state contains four signed relative errors between achieved
and requested targets, four target-to-reference descriptors, and five normalized
parameter indices. The target metrics are locked-phase DFE eye height, DFE eye
width, DFE decision margin, and CTLE power.

The reward follows the AutoCkt-style rule \citep{wang2020autockt}: only
unsatisfied relative errors contribute negative values, and a design within the
terminal tolerance receives 10. Nebula adds an explicit $-1$ result for simulator
failure instead of allowing absent output to crash the environment.

V1 proved that PPO actions could drive parameterized SPICE evaluations, but it
exposed two information problems. Reset built an observation from zero-filled
metrics rather than evaluating the starting circuit. In addition, all early
failures received the same $-1$ reward. In the fair no-warm-start real-SPICE
comparison, all 20 PPO trials failed at DC and produced identical reward, leaving
no signal to distinguish their actions.

\section{PPO v2: valid observations and informative failures}

PPO v2 was developed as an additive correction while retaining v1 state, reward,
action, and checkpoint behavior for reproducibility.

First, optional \emph{evaluated reset} simulates the initial circuit so the first
observation represents a genuine result. It costs one evaluation per episode but
removes the artificial-zero state. Second, the 28-dimensional v2 state retains
the 13 v1 entries and adds one validity flag, early-stage CTLE-power and
AC-peaking information, and a 12-way one-hot failure-stage code. The policy can
therefore distinguish setup, DC, AC, channel, and transient outcomes. Missing
receiver metrics are never silently interpreted as measured zeros.

Third, structured reward v2 gives an uninformative early failure a structural
floor of $-20$, below every valid graded design. Where genuine measurements
exist, individual target components and early-stage margins provide distance
information. A strict pass receives the terminal bonus plus a bounded overshoot
bonus capped at 2. The result separately records the legacy tolerance-based
terminal criterion and the exact strict target verdict.

V2 also added safe full checkpoints containing policy/value weights, optimizer
state, Python/NumPy/PyTorch RNG states, counters, hyperparameters, grids, target
IDs, schema versions, and Git identity. Incompatible schemas are rejected before
weights are loaded. Versioned step/update events record actions, boundary
clipping, metrics, failures, rewards, losses, entropy, and evaluation count.

The changes were tested in a 20-seed synthetic ablation from an infeasible
corner. Strict satisfaction was 0/20 for every variant at the small budget, so
the experiment does not prove circuit improvement. Under a uniform looser
criterion, corrected reset did not regress relative to v1 (3 wins, 17 ties), the
validity-aware state beat v1 in 13/20 seeds, and the full reward-v2 stack beat v1
in 15/20 seeds. Symmetric $\{-1,0,+1\}$ actions were worse in all 20 seeds and
were rejected. These are synthetic software-landscape results; no large
real-SPICE PPO v2 study was completed.

\section{PPO v3: expanded sizing and runtime contract}

PPO v3 builds on v2's failure semantics while expanding the design space from
five to eight variables. It adds matched MOS width $W$, length $L$, and integer
multiplier $M$. Both input transistors always share geometry, preserving the
matched-pair invariant. V3 has eight categorical action heads and retains the
$\{-1,0,+2\}$ grid deltas. The multiplier remains integer-valued through action
decoding, serialization, and SPICE rendering.

The frozen v3 state has 52 components: target errors and references; eight
normalized indices; diagnostic values and their validity flags; evaluator-stage
identity; channel loss at \SIlist{1.25;2.5;5}{GHz} and bulk delay; partial MOS
channel-area context; and an explicit false total-area-valid flag. Channel
features prevent evaluation on a new channel from appearing identical to the
training environment.

Reward v3 begins with reward v2 and tightens termination. Success requires a
successful evaluation, all four targets, peaking inside \SIrange{3}{12}{dB}, and
zero observed DFE errors. A weight-0.02 regularizer favors smaller matched-MOS
channel area when available. It cannot override feasibility and cannot establish
the competition's total-area requirement.

V3 checkpoints bind the policy to the ordered eight-parameter registry, exact
grids, channel checksum and port map, backend, fidelity, toolchain, and schema
versions. A five-head v1/v2 checkpoint cannot load as v3. The registry is shared
with Random Search, CEM, PVT, schematic export, and the UI.

The preserved v3 summary records a completed real-backend run with seed 101,
100 updates, and 692 evaluations. This demonstrates expanded-agent integration
with real SPICE. It does not prove v3 superior to v2, arbitrary-start convergence,
or near-optimality.

\section{Version comparison}

\begin{table}[htbp]
\centering
\caption{Technical evolution of PPO in Nebula.}
\label{tab:ppo-versions}
\begin{tabularx}{\textwidth}{@{}l c c X X@{}}
\toprule
Version & State & Heads & Main development & Evidence boundary \\
\midrule
v1 & 13 & 5 & First AutoCkt-style RL-to-SPICE baseline & End-to-end operation; flat early-failure reward exposed \\
v2 & 28 & 5 & Evaluated reset, validity/failure state, structured reward, resumable checkpoints & Multi-seed synthetic ablation; no large real-SPICE v2 proof \\
v3 & 52 & 8 & MOS $W/L/M$, channel context, stricter termination, partial-area regularizer & Real-backend artifact; no proven v2/v3 superiority \\
\bottomrule
\end{tabularx}
\end{table}

The progression is best understood as closure of the automation contract: v1
connected PPO to SPICE, v2 repaired information quality and reproducibility, and
v3 expanded physical sizing and bound downstream tools to a common identity.

\section{Random Search and CEM}

Random Search draws independent bounded candidates. CEM updates a distribution
toward an elite subset. Both use the same target and receiver evaluator as PPO.
A fair comparison must control evaluation count, time, seed, initialization,
target, channel, fidelity, and success definition. PPO is sequential and
grid-local, whereas the baselines can make global continuous draws; that
structural difference must accompany any short-horizon result.
